% why-qm.tex — Section 4: Why Quantum Mechanics?

\section{Why Quantum Mechanics?}
\label{sec:why-qm}

The experiential density $\rho$
(Sec.~\ref{subsec:experiential-measure}) identifies which structures
in $\mathcal{S}$ carry experiential weight: those with intermediate
self-modeling fidelity ($0 < I(B;\,M) < H(B)$).
Paper~5~\citep{ehrlich2026qm} answers the next question: what is the
algebraic structure of those systems?

The connection is direct.  A faithful self-model in Paper~5's sense
--- a tracking map $\varphi$ that is an order isomorphism --- implies
high mutual information $I(B;\,M)$ (the model captures all
operationally accessible structure of the body).  As long as the model
is not trivially perfect ($I < H$, i.e., the system has finite
capacity and lives in a larger environment), such a system sits in the
region where $\rho > 0$.  Therefore $\rho$ selects precisely the
structures to which Paper~5's theorem applies.  The ``area under
$\rho$'' across structure space is the totality of structures governed
by complex quantum mechanics --- the region where observers live.

\subsection{The chain in English}
\label{subsec:chain-english}

Paper~5 proves that a finite-dimensional system
admitting a faithful self-model is necessarily governed by complex
quantum mechanics.  The derivation proceeds through a chain of
implications, each either a theorem from the literature or a novel
result proved in the paper.  We summarize it here without proofs.

\textbf{Step 1: Self-modeling defines a sequential product.}  A
self-model does three things: \emph{test} an effect (probe the
system), \emph{update} the model (revise predictions), and \emph{test
again} (check the revision).  The composite operation ``test $a$, then
test $b$'' defines a \emph{sequential product} $a \cdot b$ on the
system's effect algebra.  This is not an abstract axiom imposed on
self-modeling; it is the core operation of self-modeling, constructed
via compressions and Peirce decompositions on the state space.

\textbf{Step 2: Sequential products satisfy van de Wetering's axioms.}
The sequential product inherits seven properties (S1--S7) from the
operational structure of self-modeling.  S1 (additivity) and S2
(continuity) are free.  S3 (unit law) requires the self-modeling map
$\phi$ --- without it, the unit law fails on general effects.  S4
(symmetry of orthogonality) follows from facial orthogonality in the
state space.  S5--S7 (compatibility conditions) are structural
consequences of the Peirce decomposition.  All seven are proved for
general dimension $n$ in Paper~5.

\textbf{Step 3: S1--S7 force Euclidean Jordan algebras.}  Van de
Wetering's Theorem~1~\citep{vandewetering2018}: any sequential product
space satisfying S1--S7 is a Euclidean Jordan algebra (EJA).  The
Jordan--von Neumann--Wigner classification~\citep{jvnw1934} then
restricts the state space to one of five families: $h_n(\mathbb{R})$,
$h_n(\mathbb{C})$, $h_n(\mathbb{H})$, spin factors, or the
exceptional algebra $\hthree$.

There is an independent conceptual justification for this step.  A
self-model is simultaneously a \emph{state} (it has a configuration)
and an \emph{effect evaluator} (it assigns probabilities to
configurations).  It lives in both the state cone and the effect cone.
This dual role \emph{is} self-duality.  Self-duality plus homogeneity
gives Jordan algebras via the Koecher-Vinberg
theorem~\citep{farautkoranyi1994}.  No other reconstruction program
has this natural justification for self-duality: the self-model is the
unique kind of system that IS both state and effect.

\textbf{Step 4: Faithful self-modeling forces local tomography.}  The
body-model composite $(B, M)$ admits a product-form sequential product
that inherits S1--S7 from the factors.  The correlation form
$\mathcal{B}(a, b) = \tau(a \cdot \phi^{-1}(b))$ is non-degenerate
(from the faithfulness of $\phi$ and the non-degeneracy of the trace
form).  Non-degeneracy forces the product effects to be linearly
independent, and minimality of the composite forces the composite
dimension to equal the product of the factor dimensions:
$\dim(V_{BM}) = \dim(V_B) \cdot \dim(V_M)$.  This is local
tomography: measurements on $B$ alone and $M$ alone determine the
joint state.

\textbf{Step 5: Local tomography selects complex.}  Local tomography
is the scalpel that cuts through the JvNW classification.  Real
quantum mechanics ($h_n(\mathbb{R})$) fails: $\dim = n(n+1)/2$ gives
$9 \neq 10$ for $n = 3$.  Quaternionic quantum mechanics
($h_n(\mathbb{H})$) fails: $\dim = n(2n-1)$ gives $36 \neq 28$ for
$n = 4$.  Spin factors and $\hthree$ fail to produce well-behaved
composites at all~\citep{bgw2020}.  Complex quantum mechanics
($h_n(\mathbb{C})$, $\dim = n^2$) is the unique survivor: $n^2 = n^2$
for all $n$.  The complex field is the Goldilocks choice --- real has
blind spots (no interference without complex phase), quaternionic is
redundant (tensor products misbehave).

\textbf{Step 6: Jordan + local tomography $\to$ $C^*$-algebra.}  Van
de Wetering's Theorem~3, reinforced by Barnum-Wilce and
Hanche-Olsen~\citep{hancheolsen1985}: an EJA that is locally
tomographic in composite promotes to $\Mnsa{n}$ --- the self-adjoint
part of a complex matrix algebra.  The involution (Hermitian conjugate)
emerges as the unique structure making the Jordan product
$a \circ b = \tfrac{1}{2}(ab + ba)$ compatible with the associative
matrix product.  It was never assumed.

\textbf{Step 7: $C^*$-algebra $\to$ QM.}  The remaining quantum
structure follows from established theorems:
\begin{itemize}[nosep]
\item \emph{Unitary dynamics}: automorphisms of $M_n(\mathbb{C})$ are
  inner (Skolem-Noether), hence unitary.
  Barandes~\citep{barandes2025}: unitary evolution is equivalent to
  stochastic indivisibility.
\item \emph{Born rule}: Gleason's theorem~\citep{gleason1957}
  ($n \geq 3$): the only probability measure on closed subspaces of a
  Hilbert space of dimension $\geq 3$ is the Born rule.  (For $n = 2$,
  Busch's generalization or the POVM extension applies; the chain
  does not require $n \geq 3$ as an additional assumption because the
  simple EJA condition excludes $n = 1$.)
\item \emph{Measurement}: division events in the Markov self-model
  correspond to quantum measurements via Barandes' episodic condition.
\end{itemize}

\subsection{What this means}
\label{subsec:what-this-means}

One operational premise (faithful self-modeling) plus four structural
assumptions (finite-dimensionality, faithfulness, minimal composite,
simple EJA) produces complex quantum mechanics.  The complex field is
never input, only output.  The involution is derived, not assumed.  The
Born rule is forced, not postulated.

This distinguishes the program from other reconstruction efforts.
Hardy~\citep{hardy2001} uses five axioms.  Chiribella, D'Ariano, and
Perinotti~\citep{cdp2011} use six, including purification.  Daki\'c
and Brukner~\citep{dakicbrukner2011} use four, including local
tomography as an assumption.
Masanes and M\"uller~\citep{masanesmuelleretal2013} distill to three
postulates.  Barnum, M\"uller, and
Ududec~\citep{barnummuellerududec2014} use four.  In each case,
multiple independent axioms are needed, and the complex field or a
property that selects it is typically assumed.

The present program answers a different question.  Other
reconstructions ask: ``what characterizes quantum mechanics among
generalized probabilistic theories?''  This program asks: ``\emph{why}
quantum mechanics?''  The answer: because self-modeling forces it.
Quantum mechanics is not merely characterized by operational axioms; it
is the unique theory compatible with a system that models itself.

\begin{remark}[Two formalisms]
\label{rem:two-formalisms}
Paper~1 works in an information-theoretic formalism (Markov processes,
mutual information $I(B;\,M)$).  Paper~5 works in an algebraic formalism
(order-unit spaces, sequential products).  The connection is: a faithful
tracking map $\varphi$ that is an order isomorphism (Paper~5) entails
high $I(B;\,M)$ (Paper~1), because $\varphi$ preserves all operationally
distinguishable states and their ordering.  The converse --- that high
mutual information implies an order isomorphism --- does not hold in
general and is not claimed.  The two formalisms are therefore related
asymmetrically: Paper~5's algebraic self-modeling is a \emph{sufficient}
condition for Paper~1's information-theoretic self-modeling, not an
equivalent one.  This means the algebraic chain (Steps~3--6) applies to
a subset of the systems that $\rho$ selects, namely those with faithful
(order-isomorphic) tracking.  A complete bridge theorem --- characterizing
exactly which information-theoretic self-modelers admit algebraic
self-models --- is not yet available.
\end{remark}

\subsection{Paper~4 and the Born rule}
\label{subsec:born}

Paper~4~\citep{ehrlich2026born} tested whether $\rho$ could
independently derive the Born rule via an information-theoretic
mechanism.  The result was cleanly negative: the experiential density
$\rho_Q$ is non-positive throughout all tested Lindblad trajectories.
The exchange Hamiltonian is ``too good'' at tracking --- the system
never enters the under-correlated regime where $\rho_Q > 0$.

This is not a failure of the program.  It is the program
self-correcting.  The Born rule does not need an information-theoretic
derivation because Gleason's theorem already provides it, via the
$C^*$-algebra that self-modeling forces (Steps 3--6).  Paper~4
confirms that the explanatory burden belongs on the algebraic chain,
not on $\rho$.  $\rho$'s job is simpler: select self-modelers.  The
algebra handles the rest.
